% !TeX root = ../thesis.tex
\chapter{THIẾT KẾ VÀ THỰC HIỆN PHẦN CỨNG}

% !TeX root = ../thesis.tex



\section{Phân tích kết cấu robot}

\subsection{Thiết kế thân}
Thân của robot bốn chân là bộ phận trung tâm của toàn bộ kết cấu robot. Bốn chân được liên kết với nhau thông qua thân, do đó thiết kế thân cần đáp ứng các yêu cầu sau: 

Thứ nhất, phải có độ bền đủ lớn để đảm bảo không xảy ra biến dạng, nứt gãy,… trong quá trình robot di chuyển. 

Thứ hai, để đạt được khả năng di chuyển linh hoạt và hiệu quả năng lượng cao, thân robot cần được thiết kế nhẹ nhất có thể. 

Ngoài ra, bên trong thân robot phải có không gian đủ lớn để lắp đặt các servo, pin, bo mạch điều khiển, module chuyển đổi điện áp và các linh kiện khác. 

Cuối cùng, 12 servo được lắp tại các khớp chân cần có thiết kế riêng cho đường đi dây và vị trí lắp đặt, vì chúng luôn chuyển động tương đối so với thân robot khi robot đi bộ hoặc chạy chậm. 

Dựa trên các yêu cầu trên, thân robot được thiết kế theo dạng liền khối, nhằm ngăn ngừa biến dạng hoặc gãy vỡ. 

Bốn góc của thân được thiết kế để kết nối với cụm khớp hông, đảm bảo không xảy ra giao thoa cơ khí khi khớp hông quay theo bất kỳ hướng nào. 

Ngoài ra, tại những vị trí có khả năng chịu tải thấp, cấu trúc được khoét rỗng hợp lý để giảm khối lượng. Mô hình 3D của cấu trúc thân robot bốn chân được thể hiện ở Hình \ref{fig12-c} bên dưới.

\begin{figure}[H]
    \centering
    \includegraphics[width=0.6\textwidth]{hardware/body_part_c.png}
    \caption{Thiết kế thân liền khối nhằm đạt độ bền cao và khối lượng nhẹ.}
    \label{fig12-c}
\end{figure}

\begin{figure}[H]
    \centering
    \includegraphics[width=0.6\textwidth]{hardware/body_part_b.png}
    \caption{Nắp thân.}
    \label{fig12-b}
\end{figure}

\begin{figure}[H]
    \centering
    \includegraphics[width=0.6\textwidth]{hardware/body_part_a.png}
    \caption{Thiết kế thân.}
    \label{fig12-a}
\end{figure}


\subsection{Thiết kế khớp hông}

Khớp hông của robot bốn chân được thiết kế để tạo ra hai bậc tự do, một cho chuyển động roll và một cho chuyển động pitch. Mỗi chuyển động được truyền động bởi một servo.

Do servo được sử dụng trong thiết kế là loại servo một trục, các ổ bi mặt bích 626zz được sử dụng ở phía đối diện trục servo nhằm giảm lực hướng kính tác dụng lên trục. Đầu trục và khớp hông được tích hợp trong quá trình chế tạo, điều này giúp thuận tiện cho việc lắp đặt và giảm sai số lắp ráp.

Ngoài ra, servo điều khiển chuyển động vung trước (tức là chuyển động pitch) của khớp hông được lắp bên trong giá đỡ khớp hông để tiết kiệm không gian. Mô hình 3D của cụm khớp hông được thể hiện trong Hình \ref{fig13}.

\begin{figure}[H]
    \centering
    \includegraphics[width=0.6\textwidth]{hardware/hip_joint_assembly.jpg }
    \caption{Cụm khớp hông cho chuyển động vung ngang và vung dọc.}
    \label{fig13}
\end{figure}

\subsection{Thiết kế chân}

Chân của robot bốn chân thường xuyên tiếp xúc với mặt đất để tạo lực đẩy, do đó thiết kế chân đóng vai trò quan trọng đối với khả năng di chuyển thành công của robot.

Chuyển động vung trước của khớp hông (phần đùi) sử dụng ổ bi để chịu lực dọc trục. Vì khối lượng của chân đã được bỏ qua trong các tính toán động học trình bày bên dưới, nên phần đùi và cẳng chân cần được thiết kế sao cho khối lượng và mô men quán tính nhỏ nhất có thể nhằm giảm sai số lý thuyết.

Trong thiết kế này, thay vì sử dụng kết nối servo trực tiếp tại khớp gối, một cơ cấu liên kết đầu bi (ball-head linkage) được sử dụng, trong đó cẳng chân được truyền động gián tiếp thông qua tay đòn (rocker arm) của servo. Cách bố trí này cho phép đặt servo ở phía trên của đùi, gần khớp hông, nơi mô men quán tính quay của chân tương đối nhỏ.

Mô hình lắp ráp một chân đơn được thể hiện trong Hình \ref{fig14}.

\begin{figure}[H]
    \centering
    \includegraphics[width=0.6\textwidth]{hardware/single_leg_assembly.jpg}
    \caption{Cụm lắp ráp một chân đơn.}
    \label{fig14}
\end{figure}

Dựa trên cấu trúc chân của robot bốn chân, có thể thấy rằng cẳng chân là bộ phận dễ bị gãy và uốn cong nhất vì nó phải chịu trọng lượng của toàn bộ robot cũng như các lực nén khi tiếp xúc với mặt đất. 

Ở đây, lực tiếp xúc với mặt đất được đơn giản hóa chỉ gồm lực pháp tuyến $F_N$ và lực tiếp tuyến $F_\mu$ trong mặt phẳng thẳng đứng.

Khi cẳng chân quay quanh trục khớp gối, nếu bỏ qua ma sát tại khớp và ma sát tiếp xúc với mặt đất, mô-men tải ngoài lớn nhất có thể được xác định theo công thức:

\begin{equation}
    T_{\text{max}} = \eta F_N L_3 \sin(\alpha)
\end{equation}

trong đó $F_N$ là lực pháp tuyến, $L_3$ là khoảng cách từ trục khớp gối đến điểm tiếp xúc mặt đất, $\alpha$ là góc giữa cẳng chân và phương thẳng đứng và $\eta$ là hệ số tải động.

Tổng khối lượng robot được ước tính ở giai đoạn thiết kế xấp xỉ $M = 2.5\text{ kg}$.

Giả sử bốn chân chịu đều toàn bộ trọng lượng của robot trong trạng thái tĩnh, lực pháp tuyến được tính là:

\begin{equation}
    F_N = \frac{Mg}{4}
\end{equation}

Trong phân tích bền số học dưới đây, sử dụng các giá trị: chiều dài cẳng chân $L_3 = 0.125\text{ m}$, góc nghiêng $\alpha = 45^\circ$, và hệ số tải động $\eta = 1.6$.


\section{Lựa chọn cơ cấu chấp hành}
Việc lựa chọn cơ cấu chấp hành có ảnh hưởng lớn đến hiệu năng di chuyển của robot bốn chân. Do trọng lượng bản thân của robot, cơ cấu chấp hành cần tạo ra mô-men xoắn lớn, đáp ứng nhanh và có kích thước nhỏ gọn.

Động cơ không chổi than được sử dụng rộng rãi nhờ đặc tính động học vượt trội, tuy nhiên chúng thường có giá thành cao và kích thước lớn.

Vì trọng tâm chính của nghiên cứu này là triển khai và đánh giá sơ bộ thuật toán dáng đi trên mô hình robot kích thước để bàn, nên servo được lựa chọn làm nguồn truyền động.

Được trang bị khả năng phản hồi góc quay trực tiếp theo thời gian thực cùng tỷ số truyền bánh răng lớn, loại servo này cho phép vi điều khiển dễ dàng thực hiện điều khiển vòng kín. Điều này là nền tảng thiết yếu để chạy các node tính toán động học nghịch, giúp robot tự bù trừ sai số trong các thử nghiệm dáng đi phức tạp và tự cân bằng chính xác hơn.

Mẫu servo được sử dụng là ST3215 và ST3215-HS, như thể hiện ở Hình \ref{fig16}, và các thông số kỹ thuật cụ thể được trình bày trong Bảng \ref{tab:1}.

\begin{figure}[H]
    \centering
    \includegraphics[width=0.6\textwidth]{hardware/servo_st3215.jpg}
    \caption{Cơ cấu chấp hành servo ST3215 dùng cho robot bốn chân.}
    \label{fig16}
\end{figure}

\begin{table}[htbp]
    \centering
    \caption{Thông số kỹ thuật của động cơ servo ST3215 và ST3215-HS.}
    \label{tab:1}
    \resizebox{\textwidth}{!}{% Thu nhỏ bảng vừa với chiều rộng trang giấy
        \begin{tabular}{l c c c c c c c}
            \toprule
            \textbf{Mẫu} & \textbf{Trọng lượng} & \textbf{\makecell{Kích thước \\ (L $\times$ W $\times$ H)}} & \textbf{\makecell{Điện áp \\ hoạt động}} & \textbf{\makecell{Tốc độ truyền \\ (Baudrate)}} & \textbf{\makecell{Phạm vi \\ góc quay}} & \textbf{Tốc độ tối đa} & \textbf{Lực xoắn tối đa} \\
            \midrule
            ST3215 & 69 g & \makecell{$45.2 \times 24.7 \times$ \\ $35.0$ mm} & 6.0 V -- 12.6 V & \textasciitilde 1 Mbps & \makecell{360$^\circ$ } & 0.222 s/60$^\circ$ (@12 V) & 30 kg$\cdot$cm (@12 V) \\
            \addlinespace % Thêm một khoảng trắng nhỏ giữa 2 hàng cho thoáng
            ST3215-HS & 68 g & \makecell{$45.2 \times 24.7 \times$ \\ $35.0$ mm} & 6.0 V -- 12.6 V & \textasciitilde 1 Mbps & \makecell{360$^\circ$ } & 0.094 s/60$^\circ$ (@12 V) & 20 kg$\cdot$cm (@12 V) \\
            \bottomrule
        \end{tabular}
    }
\end{table}

Để kiểm chứng mô-men xoắn định mức của servo đã chọn, tiến hành phân tích như sau. Dựa trên phân tích ở Mục 3.1.3, ảnh hưởng của mô-men xoắn lên cẳng chân là đáng kể nhất. Gọi góc giữa lực pháp tuyến tiếp xúc mặt đất $F_N$ và cẳng chân $L_3$ là $\alpha$. Như thể hiện trong Hình \ref{fig17}, mô-men do lực tiếp đất tác dụng lên khớp gối là:
\begin{equation}
    M_{\text{knee}} = \eta F_N L_3 \sin(\alpha)
\end{equation}
Mô-men này được cân bằng bởi lực truyền $F_t$ dọc theo thanh liên kết $CD$:
\begin{equation}
    M_{\text{knee}} = F_t L_{CB} \sin(\gamma)
\end{equation}
với $\gamma$ là góc giữa thanh liên kết $CD$ và cánh tay đòn $CB$. Mô-men tác dụng lên trục servo (qua cánh tay đòn $DE$) là:
\begin{equation}
    T_{\text{knee}} = F_t L_{DE} \sin(\beta)
\end{equation}
với $\beta$ là góc giữa thanh $CD$ và tay đòn $DE$. Trong thiết kế cơ cấu gần với hình bình hành, $\gamma \approx \beta$, do đó tỉ số $\sin(\beta)/\sin(\gamma) \approx 1$. Từ đó ta suy ra mô-men yêu cầu tại servo cẳng chân là:
\begin{equation}
    T_{\text{knee}} \approx \eta F_N L_3 \sin(\alpha) \frac{L_{DE}}{L_{CB}}
\end{equation}
Trong dáng đi trot, luôn có hai chân tiếp xúc mặt đất, do đó tải trọng lớn nhất tác dụng lên một chân là:
\begin{equation}
    F_N = \frac{Mg}{2}
\end{equation}
Thay các giá trị từ Bảng \ref{tab:2} với $M = 2.5\,\text{kg}$, $L_3 = 0.125\,\text{m}$, $L_{DE} = 0.024\,\text{m}$, $L_{CB} = 0.025\,\text{m}$, hệ số tải động $\eta = 1.6$, góc nghiêng $\alpha = 45^\circ$, và $g = 9.81\,\text{m/s}^2$ vào công thức, ta có $F_N = 12.26\,\text{N}$ và:
\begin{equation}
    T_{\text{knee}} \approx 1.6 \times 12.26 \times 0.125 \times \sin(45^\circ) \times \frac{0.024}{0.025} \approx 1.665\,\text{N}\cdot\text{m} \approx 17.0\,\text{kg}\cdot\text{cm}
\end{equation}

Giá trị này nhỏ hơn mô-men định mức $20\,\text{kg}\cdot\text{cm}$ của servo ST3215-HS, đảm bảo động cơ hoạt động an toàn theo thiết kế lý thuyết.

\begin{figure}[H]
    \centering
    \includegraphics[width=0.6\textwidth]{hardware/servo_torque_analysis.png}
    \caption{Phân tích mô-men xoắn của cơ cấu chấp hành servo tại cẳng chân.}
    \label{fig17}
\end{figure}

\subsection{Phân tích mô-men khớp vai (Hip Roll)}

Trong thực tế vận hành, để duy trì khả năng thăng bằng và giảm chấn, chân robot không bao giờ duỗi thẳng hoàn toàn. Thay vào đó, phần đùi ($L_2$) và cẳng chân ($L_3$) luôn tạo với nhau một góc gập nhất định như thể hiện ở Hình \ref{fig18}. Việc tính toán mô-men xoắn cần sử dụng khoảng cách thực tế từ khớp hông đến bàn chân. 

Gọi $L_{hd}$ là khoảng cách hiệu dụng từ trục khớp hông đến điểm tiếp xúc mặt đất của bàn chân. Dựa trên định lý hàm số cosin cho tam giác tạo bởi đùi, cẳng chân và góc gập gối $\beta$, ta có:
\begin{equation}
    L_{hd} = \sqrt{L_2^2 + L_3^2 - 2L_2L_3\cos(\beta)}
\end{equation}

Dựa theo Bảng \ref{tab:2} với thiết kế $L_2 = 0.106\,\text{m}$ và $L_3 = 0.125\,\text{m}$, giả sử ở tư thế đi bộ trot tiêu chuẩn robot duy trì góc khuỵu gối $\beta = 90^\circ$ nhằm tối ưu lực bám. Chiều dài hiệu dụng lúc này là:
\begin{equation}
    L_{hd} = \sqrt{0.106^2 + 0.125^2 - 2(0.106)(0.125)\cos(90^\circ)} \approx 0.164\,\text{m}
\end{equation}

Tương tự phân tích trước, tải trọng lớn nhất tác dụng lên một chân là $F_N = 12.26\,\text{N}$, hệ số tải động $\eta = 1.6$, dẫn đến lực tác dụng tổng cộng là $\eta F_N = 19.62\,\text{N}$. 

Công thức tính mô-men xoắn cực đại cho khớp vai $T_{roll}$ theo góc dang chân $\theta_1$ được xác định dựa trên khoảng lệch ngàm $L_1$ và hình chiếu ngang của chiều dài chân hiệu dụng:
\begin{equation}
    T_{roll} = \eta F_N \big(L_1 \cos(\theta_1) + L_{hd} \sin(\theta_1)\big)
\end{equation}

Thay các thông số thiết kế ($L_1 = 0.026\,\text{m}$) vào phương trình, ta thu được:
\begin{equation}
    T_{roll} = 19.62 \times \big(0.026 \cos(\theta_1) + 0.164 \sin(\theta_1)\big) \quad \text{(N}\cdot\text{m)}
\end{equation}

Dựa vào phương trình trên, có thể đánh giá mô-men xoắn trong ba trường hợp tiêu biểu. Đầu tiên, khi $\theta_1 = 0^\circ$, mô-men tác dụng lên khớp vai chỉ xuất phát từ khoảng lệch ngàm cơ khí, đạt giá trị $T_{roll} = 19.62 \times 0.026 \approx 0.51\,\text{N}\cdot\text{m} \approx 5.2\,\text{kg}\cdot\text{cm}$. Tiếp theo, khi $\theta_1 = 15^\circ$, mô-men tăng lên mức $T_{roll} \approx 1.33\,\text{N}\cdot\text{m} \approx 13.5\,\text{kg}\cdot\text{cm}$, phản ánh trạng thái robot điều chỉnh tư thế rẽ với góc dang trung bình. Cuối cùng, khi $\theta_1 = 30^\circ$, mô-men đạt giá trị lớn nhất là $T_{roll} \approx 2.05\,\text{N}\cdot\text{m} \approx 20.9\,\text{kg}\cdot\text{cm}$, tương ứng với tình huống bất lợi khi chân dang rộng.

Giá trị mô-men cực đại $20.9\,\text{kg}\cdot\text{cm}$ vẫn nhỏ hơn định mức $30\,\text{kg}\cdot\text{cm}$ của servo ST3215, do đó khớp vai hoạt động trong vùng an toàn, đảm bảo hiệu năng và không gây quá tải nhiệt.

\begin{figure}[H]
    \centering
    \includegraphics[width=0.6\textwidth]{hardware/free_body_diagram.png}
    \caption{Mô hình một chân của robot bốn chân.}
    \label{fig18}
\end{figure}

\subsection{Phân tích mô-men khớp hông (Hip Pitch)}

Tương tự quy trình đánh giá ở khớp vai, mô-men xoắn yêu cầu cực đại đối với khớp hông (Pitch) được xác định thông qua chiều dài sải bước hiệu dụng của robot. Trong mặt phẳng dọc, cánh tay đòn tạo ra mô-men phụ thuộc vào độ vươn xa của bàn chân theo phương ngang. Gọi $\alpha$ là góc nghiêng của chân (đường nối từ trục khớp hông đến bàn chân) so với phương thẳng đứng. Phương trình tổng quát tính mô-men xoắn cho khớp hông là:
\begin{equation}
    T_{pitch} = \eta F_N L_{hd} \sin(\alpha)
\end{equation}

Sử dụng khoảng cách hiệu dụng $L_{hd} = 0.164\,\text{m}$ và lực tải $\eta F_N = 19.62\,\text{N}$ đã xác định:
\begin{equation}
    T_{pitch} = 19.62 \times 0.164 \sin(\alpha) = 3.218 \sin(\alpha) \quad \text{(N}\cdot\text{m)}
\end{equation}

Phân tích quỹ đạo chuyển động cho thấy:
Phân tích quỹ đạo chuyển động cho thấy ba trạng thái chịu tải đặc trưng. Tại tư thế đứng thẳng chân ($\alpha = 0^\circ$), trọng lượng truyền dọc theo chân nên khớp hông hoàn toàn không chịu tải uốn ($T_{pitch} = 0$). Trong trường hợp khi chân bước với góc $\alpha = 15^\circ$, mô-men tại khớp hông đạt $T_{pitch} = 3.218 \sin(15^\circ) \approx 0.83\,\text{N}\cdot\text{m} \approx 8.5\,\text{kg}\cdot\text{cm}$; đây được xem là vùng tải nhẹ và rất thuận lợi cho động cơ hoạt động. Tuy nhiên, khi góc bước mở rộng lên $\alpha = 45^\circ$, mô-men tăng lên mức $T_{pitch} = 3.218 \sin(45^\circ) \approx 2.27\,\text{N}\cdot\text{m} \approx 23.2\,\text{kg}\cdot\text{cm}$, tương ứng với trạng thái vươn dài chân để leo dốc hoặc vượt chướng ngại vật.

Giá trị cực đại này vẫn nhỏ hơn giới hạn mô-men $30\,\text{kg}\cdot\text{cm}$ của servo ST3215, xác nhận rằng việc lựa chọn động cơ ST3215 cho khớp hông là hoàn toàn phù hợp và đảm bảo độ an toàn cao.

\subsection{Thiết kế tổng thể}

Thiết kế tổng thể của robot bốn chân được thể hiện trong Hình \ref{fig19}, với tổng cộng 12 bậc tự do và cấu hình chân full-elbow. Ở đây, cả bốn chân đều giống hệt nhau về cấu trúc và kích thước nhằm tạo thuận lợi cho việc điều khiển.

Khớp hông được kết nối trực tiếp với servo ở cả hai bậc tự do để đơn giản hóa cấu trúc bên trong, trong khi khớp gối được thiết kế sử dụng cơ cấu tay đòn servo kết hợp với thanh liên kết nhằm giảm quán tính của chân.

Các tham số chính của hệ thống được trình bày trong Bảng \ref{tab:2}.

\begin{figure}[H]
    \centering
    \includegraphics[width=0.6\textwidth]{hardware/preview_design.png }
    \caption{3D thiết kế robot Dog}
    \label{fig19}
\end{figure}

\begin{table}[h]
\centering
\caption{Các thông số chính của robot bốn chân}
 \label{tab:2}
\begin{tabular}{ccccccc}
\toprule
Chiều dài & Chiều rộng & Chiều cao & \makecell{Chiều dài\\hông\\$L_1$} & \makecell{Chiều dài\\đùi\\$L_2$} & \makecell{Chiều dài\\cẳng chân\\$L_3$} & \makecell{Bậc\\tự do} \\
\midrule
209 mm & 191 mm & 185 mm & 26 mm & 106 mm & 125 mm & 12 \\
\bottomrule
\end{tabular}
\end{table}

\section{Tổng quan về kiến trúc hệ thống }

Kiến trúc phần cứng của hệ thống robot bốn chân được xây dựng nhằm mục tiêu cung cấp nguồn năng lượng ổn định, đảm bảo khả năng xử lý thuật toán phức tạp theo thời gian thực và điều khiển đồng bộ nhiều cơ cấu chấp hành. Sơ đồ kết nối phần cứng tổng thể của mô hình thực nghiệm được thể hiện chi tiết trong Hình~\ref{fig30}.

\begin{figure}[H]
    \centering
    \includegraphics[width=\textwidth]{architecture/system_overview.jpeg}
    \caption{Sơ đồ kiến trúc tổng thể của hệ thống robot bốn chân.}
    \label{fig30}
\end{figure}

Hình~\ref{fig:real_robot} trình bày mô hình robot bốn chân thực tế đã được chế tạo và lắp ráp hoàn chỉnh. Robot sử dụng khung thân và các chi tiết chân được gia công bằng công nghệ in 3D FDM với vật liệu PLA, tích hợp 12 động cơ bus servo ST3215 tại các khớp, cảm biến IMU BNO055 gắn trên thân và máy tính nhúng Jetson Nano làm bộ xử lý trung tâm.

\begin{figure}[H]
    \centering
    \includegraphics[width=0.85\textwidth]{real_robot/quadrupedrobot.jpeg}
    \caption{Mô hình robot bốn chân thực tế đã được chế tạo và lắp ráp hoàn chỉnh.}
    \label{fig:real_robot}
\end{figure}

\subsection{Khối xử lý trung tâm }

Bộ xử lý trung tâm phải đáp ứng khả năng tính toán các bài toán phức tạp như động học nghịch, quy hoạch quỹ đạo và xử lý dữ liệu cân bằng theo thời gian thực. Do đó, máy tính nhúng Jetson Nano Developer Kit B01 \cite{jetson_nano} đã được lựa chọn. Vi xử lý Quad-core ARM Cortex-A57 cùng hệ điều hành Ubuntu mang lại sức mạnh vượt trội, cho phép triển khai hệ sinh thái ROS~2 và các thư viện toán học phức tạp một cách dễ dàng.

\begin{figure}[H]
    \centering
    \includegraphics[width=0.8\textwidth]{hardware/jetson_b01_specs.png} 
    \caption{Thông số kỹ thuật Jetson Nano B01}
    \label{fig34}
\end{figure}

Jetson Nano thực thi vòng lặp điều khiển chính, quản lý 12 động cơ servo và đọc dữ liệu từ module cảm biến IMU. Thiết bị thiết lập giao tiếp UART tốc độ cao với hai mạch Bus Servo Adapter (A) nhằm phân phối tín hiệu điều khiển nối tiếp đến các cơ cấu chấp hành, giúp giảm tải dòng điện và tối giản hóa sơ đồ đi dây vật lý. 

\subsection{Hệ thống Cung cấp Nguồn điện}

Trong giai đoạn phát triển và tinh chỉnh động học, hệ thống yêu cầu một nguồn điện có khả năng cung cấp dòng xả lớn và liên tục để tránh hiện tượng sụt áp khi 12 động cơ servo cùng khởi chạy. Do đó, nguồn xung tổ ong 12~V - 40~A (480~W) đã được lựa chọn. Thiết bị này đảm bảo cung cấp điện áp 12~V ổn định, tối ưu hóa hiệu suất của các servo ST3215 và hỗ trợ tốt cho quá trình thử nghiệm dài hạn trong phòng thí nghiệm.

Việc lựa chọn nguồn xung tổ ong 12~V - 40~A (480~W) được dựa trên tính toán dòng điện tiêu thụ của 12 động cơ servo ST3215:

Dòng điện không tải/hoạt động bình thường : Khoảng 300~mA - 600~mA / servo (tại 12~V).
    
Dòng điện khởi động/giữ tải lớn nhất ($I_{stall}$): Khoảng 2.7~A / servo (tại 12~V).

Trong trường hợp xấu nhất khi robot khởi động hoặc thực hiện động tác nâng toàn thân, giả sử cả 12 servo đều chịu tải lớn, tổng dòng điện tức thời ($I_{total}$) là:

$$I_{total} = N_{servo} \times I_{stall} = 12 \times 2.7 = 32.4A$$

Theo đó, tổng công suất tiêu thụ điện tối đa ($P_{total}$) của riêng hệ thống cơ cấu chấp hành là:

$$P_{total} = I_{total} \times U = 32.4 \times 12 = 388.8W$$

Để đảm bảo an toàn và tuổi thọ cho linh kiện, tránh sụt áp nguồn khi ép tải, hệ số an toàn ($k$) thường được chọn từ 1.2 đến 1.5. Ở đây nhóm chọn hệ số $k = 1.2$, ta có yêu cầu công suất nguồn cấp tối thiểu phải đạt:

$$P_{source} \ge P_{total} \times k = 388.8 \times 1.2 = 466.56W$$

Việc lựa chọn nguồn xung tổ ong 12~V - 40~A (công suất cực đại 480~W) là hoàn toàn chính xác và bám sát với tính toán lý thuyết ($466.56~W \le 480~W$). Mức công suất và dòng điện này đảm bảo khả năng đáp ứng dòng xả cao ngay cả khi robot hoạt động ở trạng thái tải nặng nhất mà không gây sụt áp hệ thống, giúp vi xử lý trung tâm Jetson Nano và các module không bị reset đột ngột.

\begin{table}[H] % <--- Chỗ này bắt buộc phải là [H]
\centering
\caption{Thông số kỹ thuật}
\renewcommand{\arraystretch}{1.3}
\begin{tabular}{p{0.3\textwidth} p{0.6\textwidth}}
\toprule
\textbf{Thuộc tính} & \textbf{Giá trị} \\
\midrule
Điện áp ngõ ra & 12~V \\
Dòng điện ngõ ra & 40-49A \\
Số điện áp ra & 1 \\
Điện áp ngõ vào & 220VAC \\
Series & Khác \\
Công suất & 400-499W \\
Cách gắn & Gắn sàn \\
Vỏ bọc & Vỏ kim loại \\
\bottomrule
\end{tabular}
\end{table}

Mức công suất 480W cung cấp một biên độ an toàn dư dả để cấp nguồn song song cho cả cụm 12 servo và máy tính nhúng Jetson Nano. Mặc dù giới hạn phạm vi di chuyển bởi dây cáp nguồn, đây là giải pháp kinh tế và đáng tin cậy nhất cho nguyên mẫu thử nghiệm tại chỗ.

\subsection{Bus Servo Adapter (A)}

Các động cơ servo ST3215 được sử dụng trong thiết kế hoạt động dựa trên chuẩn giao tiếp nối tiếp bán song công (half-duplex UART). Để máy tính nhúng Jetson Nano có thể điều khiển đồng thời cả 12 động cơ này một cách an toàn và tối ưu, hệ thống sử dụng hai mạch chuyển đổi tín hiệu Bus Servo Adapter (A). Mạch Adapter đóng vai trò trung gian, chuyển đổi tín hiệu UART tiêu chuẩn từ Jetson Nano (thông qua cổng USB) thành tín hiệu bus nối tiếp cấp cho các động cơ. 

Mặc dù về mặt lý thuyết, một mạch Bus Servo Adapter (A) có thể điều khiển tối đa 253 động cơ nối tiếp, nhưng trong thực tế, dòng điện khởi động của 12 động cơ ST3215 có thể lên đến 32.4A. Việc sử dụng song song hai mạch Adapter mang lại những lợi ích thiết thực:
Việc sử dụng song song hai mạch Adapter mang lại hai lợi ích thiết thực chính. Lợi ích thứ nhất là khả năng phân tán dòng điện hiệu quả; nhờ đó, mỗi mạch Adapter chỉ cần chịu tải dòng điện cho 6 động cơ (khoảng 16.2A tối đa), giúp giảm thiểu đáng kể rủi ro quá tải nhiệt và cháy nổ cho bản mạch in (PCB) của Adapter so với việc phải gánh toàn bộ tải. Lợi ích thứ hai là khả năng tối ưu hóa sơ đồ đi dây; thông qua việc phân chia hệ thống thành hai cụm độc lập (như nửa trái và nửa phải của robot), kết nối vật lý trở nên gọn gàng, logic và thuận tiện hơn rất nhiều cho quá trình kiểm tra cũng như bảo trì.

Thông qua các giao thức truyền thông và tập lệnh được định nghĩa sẵn, hệ thống có thể điều khiển hiệu quả các biến số như: góc quay, tốc độ di chuyển, giới hạn tải, cũng như đọc các tín hiệu phản hồi về nhiệt độ, điện áp và vị trí thực tế của từng động cơ.

\begin{figure}[H]
    \centering
    \includegraphics[width=0.6\textwidth]{hardware/bus_servo_adapter.jpg} 
    \caption{Bus Servo Adapter (A)}
    \label{fig36}
\end{figure}

\begin{table}[H]
    \centering
    \caption{Thông số kỹ thuật mạch Bus Servo Adapter (A)}
    \label{tab:adapter_specs}
    \renewcommand{\arraystretch}{1.5} % Tăng độ rộng giữa các hàng cho thoáng và đẹp giống ảnh
    \begin{tabular}{p{0.35\textwidth} p{0.5\textwidth}}
        \hline
        \textbf{Power Supply} & DC 9 $\sim$ 12.6V (ST series servo) \\ \hline
        \textbf{Power supply connector} & 5.5 $\times$ 2.1mm DC \\ \hline
        \textbf{communication interface} & UART \\ \hline
        \textbf{Dimensions} & 42.00 $\times$ 33.00mm \\ \hline
        \textbf{Mounting hole size} & 2.50mm \\ \hline
        \textbf{Mounting hole spacing} & 37.00 $\times$ 28.00mm \\ \hline
    \end{tabular}
\end{table}

\subsection{Bộ đo lường quán tính }

Module IMU (Inertial Measurement Unit) đóng vai trò cốt lõi trong hệ thống phản hồi thăng bằng của robot. Thay vì sử dụng các cảm biến quán tính thô, nghiên cứu này ứng dụng module GY-BNO055 \cite{bno055} --- một board mạch breakout nhỏ gọn được thiết kế trên nền tảng chip cảm biến BNO055 của Bosch Sensortec. Module tích hợp sẵn gia tốc kế 3 trục, con quay hồi chuyển 3 trục, từ kế 3 trục và đặc biệt là một vi xử lý ARM Cortex-M0+ 32-bit chạy thuật toán kết hợp cảm biến (Sensor Fusion) độc quyền của Bosch, cho phép xuất trực tiếp dữ liệu định hướng tuyệt đối mà không cần bộ lọc ngoại vi.

\begin{figure}[H]
\centering
\includegraphics[width=0.55\textwidth]{hardware/imu_bno055.png}
\caption{Module cảm biến 9 trục thông minh GY-BNO055 (mặt trước và mặt sau)}
\label{fig37}
\end{figure}

Module GY-BNO055 sở hữu kích thước cực kỳ nhỏ gọn (20,9~mm $\times$ 11,9~mm), rất phù hợp để gắn trực tiếp lên khung thân robot mà không chiếm nhiều không gian. Board mạch được bố trí 8 chân tín hiệu gồm: VIN (nguồn cấp), GND (mass), SCL/Rx (xung nhịp I2C hoặc nhận UART), SDA/Tx (dữ liệu I2C hoặc phát UART), ADD (chọn địa chỉ I2C), INT (ngắt ngoài), BOOT (chế độ nạp firmware) và RESET (khởi động lại phần cứng). Mạch tích hợp sẵn bộ ổn áp nội cho phép cấp nguồn linh hoạt từ 3,3V đến 5V qua chân VIN.

Về mặt kiến trúc phần cứng, nhóm đã tối ưu hóa hệ thống bằng cách loại bỏ hoàn toàn vi điều khiển trung gian như cấu trúc dùng STM32. Nền tảng phát triển giờ đây được triển khai trực tiếp trên môi trường hệ điều hành Linux của Jetson Nano. Cảm biến BNO055 được kết nối và giao tiếp thẳng với Jetson Nano thông qua bus I2C (sử dụng hai chân SCL và SDA). Việc lập trình thu thập dữ liệu được thực hiện bằng ngôn ngữ Python. Sự thay đổi này không chỉ giúp giảm thiểu đáng kể độ trễ truyền dữ liệu mà còn tối giản hóa độ phức tạp của việc đi dây phần cứng.

Khác với các cảm biến truyền thống (như MPU6050) dễ bị nhiễu từ môi trường và đòi hỏi người lập trình phải tự xây dựng các bộ lọc toán học phức tạp (như Complementary Filter hay Kalman Filter) trên máy tính chủ, BNO055 xử lý toàn bộ khối lượng công việc này ở cấp độ phần cứng nội bộ. Tính năng Sensor Fusion của BNO055 tự động tính toán và xuất trực tiếp dữ liệu định hướng dưới dạng các góc Euler (Roll, Pitch, Yaw) hoặc Quaternion với độ chính xác cao và tần số cập nhật nhanh. Do đó, bài toán của nhóm được chuyển từ việc ``viết bộ lọc nhiễu'' sang việc quản lý quy trình tự hiệu chuẩn của cảm biến.
\begin{table}[H]
\centering
\caption{Thông số kỹ thuật module cảm biến GY-BNO055}
\label{tab:bno055_specs}
\renewcommand{\arraystretch}{1.3}
\begin{tabular}{p{0.3\textwidth} p{0.55\textwidth}}
\toprule
\textbf{Thông số} & \textbf{Giá trị} \\
\midrule
Chip cảm biến & Bosch BNO055 (SiP -- System in Package) \\
Số trục đo & 9 trục (Accelerometer + Gyroscope + Magnetometer) \\
Điện áp hoạt động & 3.3V -- 5V (qua chân VIN, tích hợp bộ ổn áp nội) \\
Chuẩn giao tiếp & I2C / UART (chọn qua chân BOOT) \\
Địa chỉ I2C mặc định & 0x28 (chân ADD = LOW); 0x29 (chân ADD = HIGH) \\
Gia tốc kế (Accelerometer) & $\pm$2g / $\pm$4g / $\pm$8g / $\pm$16g \\
Con quay hồi chuyển (Gyroscope) & $\pm$125$^\circ$/s đến $\pm$2000$^\circ$/s \\
Từ kế (Magnetometer) & $\pm$1300$\mu$T (trục x,y); $\pm$2500$\mu$T (trục z) \\
Bộ xử lý tích hợp & ARM Cortex-M0+ (chạy thuật toán Sensor Fusion) \\
Dữ liệu đầu ra & Quaternion, góc Euler, vector trọng lực, gia tốc tuyến tính \\
Tần số cập nhật Fusion & Lên đến 100 Hz \\
Kích thước mạch & 20,9 $\times$ 11,9 mm \\
Số chân tín hiệu & 8 (VIN, GND, SCL/Rx, SDA/Tx, ADD, INT, BOOT, RESET) \\
Nhiệt độ hoạt động & $-40^\circ$C đến $+85^\circ$C \\
\bottomrule
\end{tabular}
\end{table}
\subsection{Mạch Giảm Áp}

Máy tính nhúng Jetson Nano yêu cầu một nguồn DC 5~V ổn định và có khả năng tiêu thụ dòng điện lên đến 4~A khi xử lý các tác vụ nặng. Để đáp ứng yêu cầu này, mạch giảm áp xung XL4015 5~A 75~W được sử dụng để chuyển đổi điện áp 12~V từ nguồn tổ ong xuống 5~V. XL4015 mang lại hiệu suất chuyển đổi cao (lên đến 96\%), dòng đầu ra tối đa 5~A, cùng các tính năng bảo vệ ngắn mạch và quá nhiệt, đảm bảo an toàn tuyệt đối cho Jetson Nano.

\begin{figure}[H]
    \centering
    \includegraphics[width=0.6\textwidth]{hardware/buck_converter.jpg} 
    \caption{Mạch giảm áp XL4015 5A 75W}
    \label{fig38}
\end{figure}

Các thông số kỹ thuật chi tiết của mạch giảm áp XL4015 được trình bày cụ thể trong Bảng \ref{tab:xl4015_specs}. Mức điện áp ngõ ra đã được chúng em tinh chỉnh qua biến trở tích hợp trên mạch để đạt chính xác mức 5.0V - 5.1V nhằm bù trừ sụt áp trên dây dẫn trước khi cấp vào Jetson Nano.

\begin{table}[H]
    \centering
    \caption{Thông số kỹ thuật mạch giảm áp XL4015 5A 75W}
    \label{tab:xl4015_specs}
    \renewcommand{\arraystretch}{1.5} 
    \begin{tabular}{p{0.3\textwidth} p{0.55\textwidth}}
        \hline
        \textbf{Điện áp ngõ vào} & 5 $\sim$ 32 VDC \\ \hline
        \textbf{Điện áp ngõ ra} & 1.25 $\sim$ 30 VDC (Tinh chỉnh về 5V cấp cho vi xử lý) \\ \hline
        \textbf{Dòng ngõ ra} & 0 $\sim$ 5A \\ \hline
        \textbf{Công suất ngõ ra} & 75W \\ \hline
        \textbf{Nhiệt độ làm việc} & -40 $\sim$ +85$^\circ$C \\ \hline
        \textbf{Tần số hoạt động} & 180KHz \\ \hline
        \textbf{Hiệu suất tối đa} & 96\% \\ \hline
        \textbf{Bảo vệ ngắn mạch} & Có (giới hạn 8A) \\ \hline
        \textbf{Bảo vệ quá nhiệt} & Có \\ \hline
    \end{tabular}
\end{table}


\subsection{Cơ cấu chấp hành: Động cơ Bus Servo ST3215}

Việc lựa chọn cơ cấu chấp hành có ảnh hưởng trực tiếp đến hiệu năng di chuyển, khả năng giữ thăng bằng và độ phức tạp của hệ thống phần cứng robot bốn chân. Trong nghiên cứu này, động cơ Bus Servo ST3215 và phiên bản tốc độ cao ST3215-HS \cite{waveshare_st3215} đã được lựa chọn làm thiết bị truyền động cho toàn bộ 12 bậc tự do của robot. 

Để làm rõ tính ưu việt và lý do lựa chọn công nghệ truyền động này cho mô hình nguyên mẫu, nhóm đã tiến hành đánh giá và so sánh dòng động cơ Bus Servo nối tiếp với hai loại cơ cấu chấp hành phổ biến khác: Động cơ Servo điều khiển bằng xung PWM truyền thống và Động cơ không chổi than (BLDC).

\begin{table}[H]
\centering
\caption{Bảng so sánh các loại cơ cấu chấp hành áp dụng cho Robot bốn chân}
\label{tab:motor_comparison}
\renewcommand{\arraystretch}{1.3}
\resizebox{\textwidth}{!}{
\begin{tabular}{|p{0.2\linewidth}|p{0.2\linewidth}|p{0.2\linewidth}|p{0.25\linewidth}|}
\hline
\textbf{Tiêu chí} & \textbf{Động cơ Servo PWM truyền thống} & \textbf{Động cơ BLDC (Tích hợp Driver)} & \textbf{Động cơ Bus Servo nối tiếp} \\ \hline
\textbf{Giao thức điều khiển} & Xung PWM (Đơn hướng) & CAN / RS485 (Song công) & UART bán song công (Half-duplex) \\ \hline
\textbf{Hệ thống dây dẫn} & Rất phức tạp (Mỗi motor 1 dây tín hiệu riêng) & Tối giản (Nối tiếp chuỗi - Daisy chain) & Tối giản (Nối tiếp chuỗi - Daisy chain) \\ \hline
\textbf{Khả năng phản hồi (Feedback)} & Không có & Vị trí, tốc độ, dòng điện, nhiệt độ (FOC) & Vị trí thực tế, tốc độ, tải, nhiệt độ, điện áp \\ \hline
\textbf{Mô-men xoắn} & Thường ở mức thấp đến trung bình & Rất cao & Cao (Tối ưu trong tầm giá) \\ \hline
\textbf{Khối lượng \& Kích thước} & Nhẹ & Nặng, cồng kềnh & Gọn nhẹ \\ \hline
\textbf{Giá thành} & Rất thấp & Rất cao & Trung bình - Hợp lý \\ \hline
\textbf{Độ phù hợp} & Thấp (Do không có phản hồi và thiếu lực) & Thấp (Dư thừa công suất, chi phí quá cao) & \textbf{Rất cao (Đạt điểm cân bằng tối ưu)} \\ \hline
\end{tabular}
}
\end{table}

Sau khi xác định kiến trúc truyền động tối ưu là sử dụng Bus Servo nối tiếp thay vì Servo PWM hay động cơ BLDC, nhóm đã tiến hành khảo sát và đánh giá các dòng Bus Servo phổ biến trên thị trường để tìm ra phương án phù hợp nhất cho hệ thống 12 bậc tự do của robot. Tuy nhiên, để có cái nhìn trực quan và toàn diện nhất, nhóm cũng đưa vào bảng so sánh một đại diện tiêu biểu của dòng Servo PWM là JX Servo PDI-HV5932MG 30KG -- loại động cơ có lực xoắn lớn, góc quay $180^{\circ}$, vốn rất phổ biến trong các dự án sinh viên để thấy rõ sự khác biệt giữa hai công nghệ.

Bốn đại diện tiêu biểu được đưa ra so sánh bao gồm: JX PDI-HV5932MG (Servo PWM), Hiwonder LX-16A, Waveshare ST3215 và Dynamixel XL430-W250-T.

\begin{table}[H]
\centering
\caption{So sánh các dòng Servo phổ biến trên thị trường}
\label{tab:bus_servo_comparison}
\renewcommand{\arraystretch}{1.3}
\resizebox{\textwidth}{!}{
\begin{tabular}{|p{0.14\linewidth}|p{0.17\linewidth}|p{0.18\linewidth}|p{0.21\linewidth}|p{0.20\linewidth}|}
\hline
\textbf{Tiêu chí} & \textbf{JX PDI-HV5932MG} & \textbf{Hiwonder LX-16A} & \textbf{Waveshare ST3215 (Được chọn)} & \textbf{Dynamixel XL430-W250-T} \\ \hline
\textbf{Giao thức} & Xung PWM ($180^{\circ}$) & TTL UART (Half-duplex) & TTL UART (Half-duplex) & TTL UART (Half-duplex) \\ \hline
\textbf{Điện áp} & 6.0~V - 8.4~V & 6.0~V - 8.4~V (Tối ưu 7.4~V) & 6.0~V - 12.6~V (Tối ưu 12~V) & 10.0~V - 14.8~V (Tối ưu 11.1~V) \\ \hline
\textbf{Lực xoắn tối đa} & 32.3 kg.cm (@8.4~V) & 17 kg.cm (@7.4~V) & 30 kg.cm / 20 kg.cm (@12~V) & $\sim$15.3 kg.cm (1.5 N.m @11.1~V) \\ \hline
\textbf{Cảm biến vị trí} & Biến trở (Potentiometer) & Biến trở (Potentiometer) & Từ tính (Magnetic Encoder $360^{\circ}$) & Từ tính (Absolute Encoder) \\ \hline
\textbf{Độ phân giải góc} & Phụ thuộc vi điều khiển & 0.24$^{\circ}$ & 0.088$^{\circ}$ & 0.088$^{\circ}$ \\ \hline
\textbf{Tốc độ truyền} & 50~Hz - 330~Hz (PWM) & 115200 bps & Lên đến 1~Mbps & Lên đến 4.5~Mbps \\ \hline
\textbf{Bộ điều khiển} & Cơ bản (P) & Cơ bản (P) & Nâng cao (PI) & Cao cấp (PID vị trí, vận tốc, dòng) \\ \hline
\textbf{Giá thành} & Trung bình ($\sim$500.000 VNĐ) & Thấp ($\sim$300.000 VNĐ) & Trung bình ($\sim$500.000 VNĐ) & Rất cao ($\sim$1.500.000 VNĐ) \\ \hline
\end{tabular}
}
\end{table}

Dựa trên các thông số phân tích từ Bảng \ref{tab:bus_servo_comparison}, Waveshare ST3215 được lựa chọn là giải pháp tối ưu cho hệ thống cơ cấu chấp hành vì những lý do sau:

Lựa chọn ST3215 giúp khắc phục triệt để các nhược điểm của những dòng động cơ khác. Cụ thể, đối với Servo PWM như JX PDI-HV5932MG, mặc dù có lực xoắn rất ấn tượng (lên đến 32.3 kg.cm ở mức 8.4V), góc quay giới hạn $180^{\circ}$ và cùng phân khúc giá với ST3215, nhưng việc sử dụng giao thức PWM đơn hướng khiến nó hoàn toàn không có khả năng phản hồi trạng thái (góc quay thực tế, dòng điện, nhiệt độ). Điều này tạo ra cản trở lớn đối với bài toán điều khiển thăng bằng vòng kín của robot, đồng thời sơ đồ đi dây 12 tín hiệu PWM rời rạc rất dễ sinh nhiễu và đứt ngầm. So với Hiwonder LX-16A có lợi thế giá thành thấp, ST3215 vượt trội hơn hẳn về lực xoắn (LX-16A chỉ đạt 17 kg.cm) cũng như độ phân giải (cảm biến vị trí biến trở của LX-16A có độ phân giải thấp $0.24^{\circ}$ và mau mòn, với bộ điều khiển P gây sai số tĩnh lớn). So với Dynamixel XL430-W250-T mang tính năng cao cấp nhưng giá thành đắt đỏ (gấp 3 lần) và lực xoắn thấp hơn ($\sim$15.3 kg.cm), ST3215 mang lại điểm cân bằng hoàn hảo. Lực xoắn cao của ST3215 cung cấp dư dả công suất để duy trì thăng bằng; cảm biến từ tính với độ phân giải cực cao ($0.088^{\circ}$) phản hồi góc chính xác tuyệt đối mà không bị hao mòn; tốc độ truyền tải 1~Mbps và bộ điều khiển PI tích hợp đáp ứng xuất sắc vòng lặp điều khiển thời gian thực. Đặc biệt, hệ thống còn tận dụng phân bổ chiến lược giữa hai phiên bản ST3215 và ST3215-HS để tối ưu hóa đặc tính động lực học của từng khớp. Tại khớp vai và khớp hông (Hip Roll, Hip Pitch) --- nơi phải chịu phần lớn tải trọng và tạo lực đẩy chính --- phiên bản tiêu chuẩn ST3215 (30 kg.cm) được sử dụng để duy trì hệ số an toàn cao, giúp động cơ bền bỉ và không sinh nhiệt quá mức. Ngược lại, tại khớp gối (Knee), do yêu cầu lực uốn thấp hơn nhưng lại đòi hỏi tốc độ gập duỗi cực nhanh trong pha vung chân, phiên bản tốc độ cao ST3215-HS (20 kg.cm) đã được lựa chọn; tốc độ phản hồi đạt $0.094\,\text{s}/60^{\circ}$ của nó giúp robot di chuyển linh hoạt và lấy lại thăng bằng chớp nhoáng khi gặp nhiễu.

\subsection{Hệ thống Phân phối Năng lượng và Lựa chọn Dây dẫn}

Để đảm bảo hệ thống 12 động cơ hoạt động ổn định ở dòng tải cao mà không bị sụt áp, việc thiết kế sơ đồ đi dây và lựa chọn tiết diện dây dẫn đóng vai trò cực kỳ quan trọng. Dựa trên phân tích dòng điện tổng ($I_{total} \approx 32.4A$) từ nguồn tổ ong, hệ thống dây dẫn được chia thành 3 cấp độ với các tiêu chuẩn AWG (American Wire Gauge) khác nhau, được phân phối qua các khối đầu nối nguồn (Terminal Block).

\begin{figure}[H]
    \centering
    \includegraphics[width=1.0\textwidth]{hardware/system_wiring_diagram.jpg} 
    \caption{Sơ đồ phân phối điện áp và tín hiệu tổng thể của robot}
    \label{fig44}
\end{figure}

Dựa trên thông số kỹ thuật (Datasheet) của các loại dây lõi đồng, nhóm đã đưa ra thiết kế đi dây như sau:

Đối với trục nguồn chính (từ nguồn tổ ong đến khối đấu nối), cáp AWG 12 được sử dụng. Loại cáp này đóng vai trò truyền tải toàn bộ dòng điện từ nguồn xung 12~V - 40~A đến các khối chia nguồn trung tâm. Với đường kính lõi khoảng 2.05 mm, tiết diện 3.31 mm$^2$, khả năng chịu dòng trong không khí mở đạt 30A - 40A cùng điện trở cực thấp (khoảng 5.2 $\Omega$/km), dây AWG 12 hoàn toàn đáp ứng được dòng xả tức thời 32.4A của 12 servo mà không sinh nhiệt hay gây sụt áp đáng kể.

Đối với nhánh nguồn phụ (từ khối đấu nối đến servo đầu tiên), cáp AWG 20 được áp dụng nhằm cấp nguồn trực tiếp cho chuỗi 3 servo của mỗi chân. Với đường kính lõi 0.81 mm và tiết diện 0.52 mm$^2$, loại dây này chịu được dòng định mức khoảng 5A (tối đa 11A). Vì mỗi nhánh 3 servo ST3215 có dòng khởi động tổng cộng khoảng 8.1A, dây AWG 20 cung cấp sự cân bằng hoàn hảo giữa khả năng tải dòng an toàn và độ mềm dẻo cần thiết để luồn lách linh hoạt qua các khớp cơ khí.

Cuối cùng, đối với bus tín hiệu truyền thông (từ Adapter đến servo), cáp AWG 22 được lựa chọn để truyền tín hiệu UART (Data) và nối chung chân Mass (GND) nhằm đồng bộ điện áp logic. Với đường kính lõi 0.64 mm, do đặc thù đường truyền dữ liệu chỉ tiêu thụ dòng điện vài mA, việc sử dụng dây AWG 22 giúp cáp tín hiệu trở nên nhỏ gọn, hạn chế nhiễu chéo, thuận tiện cho việc bấm cos và đi dây nối tiếp (daisy-chain) qua từng servo mà hoàn toàn không cản trở chuyển động của các khớp chân.

Việc thiết kế phân nhánh nguồn thông qua Terminal Block kết hợp với việc tính toán và lựa chọn đúng cỡ dây AWG từ datasheet giúp giải quyết triệt để vấn đề sụt áp trên đường dây (Voltage Drop), đảm bảo 12 động cơ ở cuối chuỗi vẫn nhận đủ điện áp 12~V để sinh ra mô-men xoắn tối đa.

\section{Các Giao thức Truyền thông của Hệ thống}

Hệ thống điều khiển sử dụng hai giao thức truyền thông chính: I2C (tần số 400 kHz) để thu thập dữ liệu tư thế từ cảm biến IMU BNO055 và UART (tốc độ 1 Mbps) để điều khiển 12 cơ cấu chấp hành ST3215 thông qua mạch Bus Servo Adapter (A). Việc sử dụng UART bán song công với tốc độ truyền tải cao trên hai bus vật lý riêng biệt giúp đảm bảo hệ thống duy trì được thời gian phản hồi theo yêu cầu của vòng lặp điều khiển thời gian thực.

\subsection{Phân tích thời gian truyền thông UART}

Việc phân tích thời gian truyền thông trên bus UART là cần thiết để đảm bảo tốc độ cập nhật lệnh servo đáp ứng yêu cầu thời gian thực của vòng lặp điều khiển. Với baudrate 1 Mbps và giao thức bán song công (half-duplex), mỗi giao dịch bao gồm:

Trong một giao dịch hoàn chỉnh, gói lệnh ghi (Write) bao gồm tổng cộng 10 byte, được cấu thành từ header (2 byte), ID servo (1 byte), chiều dài (1 byte), lệnh (1 byte), dữ liệu vị trí và tốc độ (4 byte), cùng checksum (1 byte). Tương ứng, gói phản hồi (Response) gồm tổng cộng 8 byte, phân bổ cho header (2 byte), ID servo (1 byte), chiều dài (1 byte), mã lỗi (1 byte), dữ liệu (2 byte) và checksum (1 byte).

Với cấu hình 8N1 (8 bit dữ liệu, không parity, 1 stop bit), mỗi byte cần 10 bit truyền. Thời gian cho một giao dịch hoàn chỉnh (ghi + đọc) với một servo:

\begin{equation}
    t_{servo} = \frac{(10 + 8) \times 10}{10^6} = 180 \; \mu\text{s}
\end{equation}

Mỗi bus driver quản lý 6 servo, thời gian tổng cho một frame lệnh:

\begin{equation}
    t_{bus} = 6 \times t_{servo} = 6 \times 180 = 1080 \; \mu\text{s} \approx 1{,}1 \text{ ms}
\end{equation}

Hai bus hoạt động song song trên hai cổng USB riêng biệt, nên tổng thời gian truyền thông cho 12 servo vẫn xấp xỉ $1{,}1$ ms --- nhỏ hơn nhiều so với chu kỳ điều khiển $\Delta T = 7$ ms. Điều này xác nhận băng thông truyền thông đủ cho vận hành thời gian thực.

\subsection{Ánh xạ ID servo và phân bổ driver}

Bảng~\ref{tab:servo_mapping} trình bày cấu hình ánh xạ ID servo và phân bổ cho hai driver UART.

\begin{table}[H]
\centering
\caption{Ánh xạ ID servo theo driver và chân robot}
\label{tab:servo_mapping}
\renewcommand{\arraystretch}{1.3}
\begin{tabular}{|l|c|c|c|c|}
\hline
\textbf{Chân} & \textbf{Driver} & \textbf{Hip Roll} & \textbf{Hip Pitch} & \textbf{Knee} \\ \hline
Left-Front (LF) & A & ID 7 & ID 8 & ID 9 \\ \hline
Right-Behind (RB) & A & ID 10 & ID 11 & ID 12 \\ \hline
Right-Front (RF) & B & ID 1 & ID 2 & ID 3 \\ \hline
Left-Behind (LB) & B & ID 4 & ID 5 & ID 6 \\ \hline
\end{tabular}
\end{table}

Chiến lược phân bổ ghép hai chân \textit{chéo nhau} trên cùng một driver (LF+RB trên Driver A, RF+LB trên Driver B) có ý nghĩa quan trọng: trong dáng đi Trot, hai chân chéo luôn ở cùng pha (vung chân hoặc chống chân). Khi một cặp chéo đang ở pha vung chân (cần cập nhật vị trí nhanh), cặp còn lại ở pha chống chân (ít thay đổi), giúp cân bằng tải truyền thông giữa hai bus.

